\documentclass{article}
\usepackage[zh]{homework}

\config{分析\,-\,0}{2026 春季}{11}

\begin{document}

\begin{problem}
设 $f \in C^3([-1,1])$, $f(0)=0$, $f(1)=1$, $f'(0)=1$, $f''(0)=0$. 求
\[
\int_0^1 (1-t)^2 f'''(t)\d t.
\]
\end{problem}

\begin{solution}
  根据积分余项的泰勒公式, 有
  \[ f(1)
  =f(0)+\frac{1}{1!}f'(0)+\frac{1^2}{2!}f''(0)+\int_{0}^{1}\frac{(1-t)^2}{2!}f'''(t)\d t. \]
  结合题目条件, 可得
  \[ \int_0^1 (1-t)^2 f'''(t)\d t=2(1-1)=0. \]
\end{solution}



\begin{problem}
判断下列反常积分的敛散性, 并说明理由.
\begin{enumerate}[label=(\alph*)]
\item $\int_2^\infty \frac{\d x}{x(\ln x)^p}$, $(p>0)$.
\item $\int_1^\infty \frac{x\ln x}{(1+x^2)^2}\d x$.
\item $\int_0^\infty \frac{\ln x}{x^p}\d x$, $(p>0)$.
\item $\int_0^{\pi/2} \frac{\d x}{\sin^p x}$, $(p>0)$.
\item $\int_0^1 \frac{\d x}{x^p - x^q}$, $(q>p>0)$.
\end{enumerate}
\end{problem}

\begin{solution}
  \begin{enumerate}[label=(\alph*)]
    \item 注意到
    \[ \int \frac{\d x}{x(\ln x)^p} = \int (\ln x)^{-p}\d \ln x
    =\begin{cases}
    \frac{1}{1-p}(\ln x)^{1-p}+C, & p\neq 1, \\
    \ln(\ln x)+C, & p=1.
    \end{cases} \]
    因此, 若 \( p>1 \), 则 \( \lim_{x\to\infty} \frac{1}{1-p}(\ln x)^{1-p} = 0 \), 积分收敛; 若 \( p\leq 1 \), 则积分发散.
    \item 注意到
    \begin{align*}
      \int \frac{x\ln x}{(1+x^2)^2}\d x
      &= \frac{1}{4} \int \frac{\ln(x^2)}{(1+x^2)^2}\d(x^2)
      = \frac{1}{4} \int \frac{\ln t}{(1+t)^2}\d t \qquad (t=x^2) \\
      &= -\frac{1}{4} \int \ln t\,\d\!\left(\frac{1}{1+t}\right)
      = -\frac{1}{4} \frac{\ln t}{1+t} + \frac{1}{4} \int \frac{1}{t(1+t)}\d t \\
      &= \frac{1}{4} \ln \frac{t}{1+t} - \frac{1}{4} \frac{\ln t}{1+t} + C.
    \end{align*}
    因此
    \[
      \int_1^\infty \frac{x\ln x}{(1+x^2)^2}\d x
      = \lim_{b\to\infty} \left[\frac{1}{4} \ln \frac{t}{1+t} - \frac{1}{4} \frac{\ln t}{1+t}\right]_{t=1}^{t=b^2}
      = \frac{\ln 2}{4}.
    \]
    因此该反常积分收敛.
    \item 若 \( p\neq 1 \), 则
    \begin{align*}
      \int \frac{\ln x}{x^{p}}\d x
      &=\int \frac{t}{e^{p t}}e^{t}\d t\quad (t=\ln x) \\
      &=\int te^{(1-p)t}\d t
      =\frac{e^{(1-p)t}}{(1-p)^2}((1-p)t-1)\\
      &=\frac{x^{1-p}\left( (1-p)\ln x-1 \right)}{(1-p)^2}
      :=F(x).
    \end{align*}

    若 \( 1-p>0 \), 则 \( \lim_{x\to+\infty}F(x)=\infty \), 积分发散; 若 \( 1-p<0 \), 则 \( \lim_{x\to 0^+}F(x)=+\infty \), 积分发散.

    若 \( p=1 \), 则
    \[
      \int \frac{\ln x}{x}\d x = \frac{1}{2}(\ln x)^2+C.
    \]
    因此积分发散.
    \item 若 \( p\ge 1 \), 则 \( 1/\sin^{p}(x)\ge1/\sin x \), 而 
    \[ \int_{0}^{\pi/2}\frac{\d x}{\sin x}=\ln\left( \tan \frac{x}{2} \right)\Big|_{0}^{\pi/2} = +\infty. \]
    因此积分发散.

    若 \( p<1 \), 由于对于任意的 \( x\in[0,\pi/2] \) 均有 \( \sin x\geq2x/\pi \), 因此
    \[ \int_{0}^{\pi/2}\frac{\d x}{\sin^p x}\geq \int_{0}^{\pi/2}\frac{\d x}{(2x/\pi)^p} = \left(\frac{\pi}{2}\right)^p \int_{0}^{\pi/2} x^{-p} \d x. \]
    当 \( p<1 \) 时, 后者收敛, 因此原积分收敛.

    \item 注意到, 对于 \( x\in[0,1] \), 有
    \[ x^{p}-x^q=x^p(1-x^{p-q})\leq 1-x^{p-q} \leq \max\{1,p-q\}(1-x). \]
    记 \( M=\max\{1,p-q\} \), 则
    \[ \int_0^1 \frac{\d x}{x^p - x^q} \geq \int_0^1 \frac{\d x}{M(1-x)} = +\infty. \]
    因此积分发散.
  \end{enumerate}
\end{solution}



\begin{problem}
利用 Cauchy 积分判别法判断下列级数的敛散性, 并给出理由.
\begin{enumerate}[label=(\alph*)]
\item $\sum_{n=2}^\infty \frac{1}{n(\ln n)^p}$, $(p>0)$.
\item $\sum_{n=3}^\infty \frac{1}{n\ln n\ln(\ln n)}$.
\end{enumerate}
\end{problem}

\begin{solution}
  (a) 这个问题与问题 2(a) 基本上等价, 只需补充以下事实: 对于足够大的 \( x \), 函数 \( y=1/(x(\ln x)^p) \) 单调递减, 因此 \( \sum_{n=2}^{+\infty}1/(n(\ln n)^p) \) 与 \( \int_{2}^{+\infty}1/(x(\ln x)^p)\d x \) 同时收敛或同时发散.

  (b) 注意到
  \begin{align*}
    \int \frac{1}{x\ln x\ln(\ln x)}\d x
    &=\int \frac{1}{t\ln t}\d t \qquad (t=\ln x) \\
    &=\ln \ln t + C = \ln(\ln(\ln x)) + C.
  \end{align*}
  再注意到, 对于充分大的 \( x \), 函数 \( y=1/(x\ln x\ln(\ln x)) \) 单调递减, 因此 \( \sum_{n=3}^{+\infty}1/(n\ln n\ln(\ln n)) \) 与 \( \int_{3}^{+\infty}1/(x\ln x\ln(\ln x))\d x \) 同时收敛或同时发散. 由于后者发散, 因此前者也发散.
  
\end{solution}



\begin{problem}
\begin{enumerate}[label=(\alph*)]
\item 分部积分有时能用于无穷限反常积分. 列出适当的假设, 编成定理, 并加以证明.
\item 证明
\[
\int_0^\infty \frac{\cos x}{1+x}\d x = \int_0^\infty \frac{\sin x}{(1+x)^2}\d x.
\]
证明这两个积分中有一个绝对收敛, 另一个则不然.
\end{enumerate}
\end{problem}

\begin{solution}
  (a) 定理: 设 \( f,g\in C^1([a,+\infty)) \), 且 \( \lim_{x\to+\infty}f(x)g(x) \) 存在, 且 \( \int_{a}^{+\infty}u'(x)v(x)\d x \) 存在. 则
  \[ \int_{a}^{+\infty}u(x)v'(x)\d x = \left[ \lim_{x\to+\infty}f(x)g(x) - f(a)g(a) \right] - \int_{a}^{+\infty}u'(x)v(x)\d x. \]

  这是因为, 对于任意 \( A>a \), 有
  \[ \int_{a}^{A}u(x)v'(x)\d x = \left[ f(A)g(A) - f(a)g(a)  \right]- \int_{a}^{A}u'(x)v(x)\d x. \]
  取 \( A\to+\infty \), 由于已知等式右侧存在, 因此等式左侧的反常积分也存在, 且满足定理中的等式.

  (b) 取 \( u(x)=\sin x \), \( v(x)=(1+x)^{-1} \). 则由于 \( \lim_{x\to+\infty}u(x)v(x)=0 \), 并且 \( u(0)v(0)=0 \), 因此根据 (a) 中的定理, 可得
  \[ \int_0^\infty \frac{\cos x}{1+x}\d x = \int_0^\infty \frac{\sin x}{(1+x)^2}\d x. \]

  下面来证明: \( \int_0^\infty \frac{\sin x}{(1+x)^2}\d x \) 绝对收敛, 而 \( \int_0^\infty \frac{\cos x}{1+x}\d x \) 不绝对收敛. 注意到
  \[ \int_{0}^{+\infty}\abs{\frac{\sin x}{(1+x)^2}}\d x\le \int_{0}^{+\infty}\frac{1}{(1+x)^2}\d x=1, \]
  因此第一个积分绝对收敛. 而
  \begin{align*}
    \int_{0}^{+\infty}\abs{\frac{\cos x}{1+x}}\d x
    &>\sum_{n=1}^{+\infty}\int_{(n-1/2)\pi}^{(n+1/2)\pi}\frac{\abs{\cos x}}{1+x}\d x
    >\sum_{n=1}^{+\infty}\frac{1}{1+(n+1/2)\pi}\int_0^{\pi}\sin x\d x\\
    &=\sum_{n=1}^{+\infty}\frac{2}{1+(n+1/2)\pi}=+\infty.
  \end{align*}
  因此第二个积分不绝对收敛.
\end{solution}



\begin{problem}
对 $1<s<\infty$, 定义
\[
\zeta(s)=\sum_{n=1}^\infty \frac{1}{n^s},
\]
(这是 Riemann 的 $\zeta$ 函数, 在研究质数的分布时极为重要.) 求证:
\begin{enumerate}[label=(\alph*)]
\item $\zeta(s)= s\int_1^\infty \frac{\lfloor x\rfloor}{x^{s+1}}\d x$, 其中 $\lfloor x\rfloor$ 表示不大于 $x$ 的最大整数.
\item $\zeta(s)= \frac{s}{s-1} - s\int_1^\infty \frac{x-\lfloor x\rfloor}{x^{s+1}}\d x$, 并证明该积分对一切 $s>0$ 收敛.
\end{enumerate}
提示: 为了证明 (a), 可以求 $[1,N]$ 上的积分与定义 $\zeta(s)$ 的级数的第 $N$ 个部分和之差.
\end{problem}

\begin{solution}
  (a) 设 \( N\in\mathbb{N} \), 则
  \begin{align*}
    s \int_{1}^{+\infty}\frac{\fl{x}}{x^{s+1}}\d x
    &= \lim_{N\to\infty} \sum_{k=1}^{N} \int_{k}^{k+1} \frac{k s}{x^{s+1}}\d x
    = \lim_{N\to\infty} \sum_{k=1}^{N} k\left( \frac{1}{k^s}-\frac{1}{(k+1)^s} \right)\\
    &= \lim_{N\to\infty} \left( -\frac{N}{(N+1)^s} + \sum_{k=1}^{N}\frac{1}{k^s} \right) 
    = \sum_{n=1}^{+\infty}\frac{1}{n^s}=\zeta(s).
  \end{align*}

  (b) 注意到
  \[ \int_{1}^{+\infty}\frac{x}{x^{s+1}}\d x
  \int_{1}^{+\infty}x^{-s}\d x
  = \frac{x^{1-s}}{1-s} \Big|_{1}^{+\infty}
  = \frac{1}{s-1}. \]
  因此 \( \int_1^\infty \frac{x-\lfloor x\rfloor}{x^{s+1}}\d x \) 收敛, 并且
  \[  \frac{s}{s-1} - s\int_1^\infty \frac{x-\lfloor x\rfloor}{x^{s+1}}\d x
  =\frac{s}{s-1}-\frac{s}{s-1} + \zeta(s)=\zeta(s). \]
\end{solution}



\begin{problem}
考虑曲线 $\gamma:[a,b]\to\R^k$. 设 $\gamma\in C^1([a,b])$, 证明 $\gamma$ 可求长, 并且
\[
l(\gamma)=\int_a^b \abs{\gamma'(t)}\d t.
\]
\end{problem}

\begin{proof}
  沿用教材中的记号. 一方面, 对于任意的划分 \( P \), 均有
  \begin{align*}
    \Lambda(P,\gamma)
    &=\sum_{k=1}^{n} \abs{\gamma(t_k)-\gamma(t_{k-1})}
    =\sum_{k=1}^{n} \abs{\int_{t_{k-1}}^{t_k} \gamma'(t)\d t}
    \leq \sum_{k=1}^{n} \int_{t_{k-1}}^{t_k} \abs{\gamma'(t)}\d t
    =\int_a^b \abs{\gamma'(t)}\d t.
  \end{align*}
  因此, \( l(\gamma) \) 存在, 并且
  \[ l(\gamma) = \int_a^b \abs{\gamma'(t)}\d t. \]

  另一方面, 对于任意的正实数 \( \varepsilon \), 由于 \( \gamma' \) 是定义在闭区间上的连续函数, 故一致连续. 即, 存在正实数 \( \delta \), 使得对于任意的 \( x,y\in[a,b] \) 满足 \( \abs{x-y}<\delta \) 时, 有 \( \abs{\gamma'(x)-\gamma'(y)}<\varepsilon \). 对于任意的划分 \( P \) 满足 \( \|P\|<\delta \), 则对于任意的 \( k\in\{1,2,\dots,n\} \) 和任意的 \( t\in[t_{k-1},t_k] \), 有
  \begin{align*}
    \abs{\gamma(t_{k})-\gamma(t_{k-1})}
    &=\abs{\int_{t_{k-1}}^{t_k} \gamma'(t)\d t}
    =\abs{\int_{t_{k-1}}^{t_k} (\gamma'(t)-\gamma'(t_{k}))\d t + \Delta t_k \gamma'(t_{k})} \\
    &\geq \Delta t_k \abs{\gamma'(t_{k})} - \int_{t_{k-1}}^{t_k} \abs{\gamma'(t)-\gamma'(t_{k})}\d t \\
    &\geq \int_{t_{k-1}}^{t_k} \abs{\gamma'(t)} \d t + \int_{t_{k-1}}^{t_k}\left( \abs{\gamma'(t_k)}-\abs{\gamma'(t)} \right) \d t - \Delta t_k \varepsilon \\
    &\geq \int_{t_{k-1}}^{t_k} \abs{\gamma'(t)} \d t -2 \Delta t_k \varepsilon.
  \end{align*}
  因此
  \[ \Lambda(P,\gamma)
  = \sum_{k=1}^{n}\abs{\gamma(t_k)-\gamma(t_{k-1})}
  \geq \sum_{k=1}^{n}\left( \int_{t_{k-1}}^{t_k} \abs{\gamma'(t)} \d t -2 \Delta t_k \varepsilon \right)
  =\int_{a}^{b} \abs{\gamma'(t)}\d t - 2\varepsilon(b-a).
  \]
  由于 \( \varepsilon \) 是任意的, 故根据定义 \( l(\gamma)=\sup_{P}\Lambda(P,\gamma) \), 有
  \[ l(\gamma) = \int_a^b \abs{\gamma'(t)}\d t. \]
\end{proof}



\begin{problem}
设复平面里的曲线 $\gamma_1,\gamma_2,\gamma_3$ 是由
\[
\gamma_1(t)=e^{\i t},\quad \gamma_2(t)=e^{2\i t},\quad \gamma_3(t)=e^{2\pi \i t\sin(1/t)}
\]
定义在 $[0,2\pi]$ 之上的, 令 $\gamma_3(0)=1$.
\begin{enumerate}[label=(\alph*)]
\item 证明这三条曲线的值域相同.
\item 证明 $\gamma_1$ 与 $\gamma_2$ 为可求长, $\ell(\gamma_1)=2\pi$, $\ell(\gamma_2)=4\pi$, 而 $\gamma_3$ 不可求长.
\end{enumerate}
\end{problem}

\begin{solution}
  \begin{enumerate}[label=(\alph*)]
    \item 对于 \( \gamma_1,\gamma_2 \), 显然有 \( \abs{\gamma_1(t)}=\abs{\gamma_2(t)}=1 \), 因此它们的值域均为单位圆 \(\{z\in\mathbb{C}:\abs{z}=1\}\).

    对于 \( \gamma_3 \), 令
    \[ \varphi(t)=2\pi t\sin(1/t)\ (t>0),\qquad \varphi(0)=0. \]
    则 \( \gamma_3(t)=e^{\i\varphi(t)} \), 且 \( \varphi \) 在 \([0,2\pi]\) 上连续, 因而 \( \varphi([0,2\pi]) \) 为一闭区间 \([m,M]\). 注意到
    \[ \varphi\left(\frac{2}{3\pi}\right)=2\pi\cdot\frac{2}{3\pi}\sin\left(\frac{3\pi}{2}\right)=-\frac{4}{3}, \]
    且由 \( \sin x\ge x-\frac{x^3}{6} \) 得
    \[ \varphi(2\pi)=4\pi^2\sin\left(\frac{1}{2\pi}\right)
    \ge 4\pi^2\left(\frac{1}{2\pi}-\frac{1}{48\pi^3}\right)=2\pi-\frac{1}{12\pi}. \]
    因此 \( M-m\ge 2\pi-\frac{1}{12\pi}+\frac{4}{3}>2\pi \). 由于指数函数 \( e^{\i x} \) 以 \( 2\pi \) 为周期, 故 \( e^{\i\varphi([0,2\pi])} \) 覆盖整个单位圆. 于是三条曲线的值域相同.

    \item \( \gamma_1'(t)=\i e^{\i t} \), \( \gamma_2'(t)=2\i e^{2\i t} \), 因此
    \[ \ell(\gamma_1)=\int_0^{2\pi}\abs{\gamma_1'(t)}\d t=\int_0^{2\pi}1\d t=2\pi,\qquad
    \ell(\gamma_2)=\int_0^{2\pi}\abs{\gamma_2'(t)}\d t=\int_0^{2\pi}2\d t=4\pi. \]

    对于 \( \gamma_3 \), 当 \( t>0 \) 时可导且
    \[ \gamma_3'(t)=2\pi\i e^{2\pi\i t\sin(1/t)}\left(\sin(1/t)-\frac{\cos(1/t)}{t}\right). \]
    下面来证明 \( \gamma_3 \) 不可求长, 即证明
    \[ \int_{0}^{2\pi}\abs{\gamma_3'(t)}\d t \]
    发散.

    对任意 \( a\in(0,2\pi] \), \( \gamma_3 \) 在 \([a,2\pi]\) 上为 \( C^1 \) 曲线, 因此
    \[ \ell\bigl(\gamma_3|_{[a,2\pi]}\bigr)
    =\int_a^{2\pi}2\pi\left\lvert\sin(1/t)-\frac{\cos(1/t)}{t}\right\rvert\d t. \]
    取 \( t_k=1/(2k\pi) \), 并令
    \[ I_k=\left[\frac{1}{2k\pi+\pi/6},\ \frac{1}{2k\pi-\pi/6}\right], \qquad k\ge 2. \]
    则当 \( t\in I_k \) 时有 \( \abs{\cos(1/t)}\ge 1/2 \), 且 \( \abs{\sin(1/t)}\le 1 \), 从而对充分大的 \( k \) 有
    \[ \left\lvert\sin(1/t)-\frac{\cos(1/t)}{t}\right\rvert\ge \frac{1}{2t}-1\ge \frac{1}{4t}. \]
    因此
    \[ \ell\bigl(\gamma_3|_{[a,2\pi]}\bigr)
    \ge \sum_{k:\,I_k\subset[a,2\pi]}\int_{I_k}\frac{\pi}{2}\frac{\d t}{t}
    =\frac{\pi}{2}\sum_{k:\,I_k\subset[a,2\pi]}\ln\frac{2k\pi+\pi/6}{2k\pi-\pi/6}, \]
    而 \( \ln\frac{2k\pi+\pi/6}{2k\pi-\pi/6}\sim c/k \) (\( c>0 \)), 故上式随 \( a\to 0^+ \) 发散.
    于是 \( \ell(\gamma_3)=+\infty \), \( \gamma_3 \) 不可求长.
  \end{enumerate}
\end{solution}



\end{document}